\begin{houseviewbox}[结论：减持 / 避免新增仓位；12 个月目标价 17.0 元]
截至 2026 年 7 月 23 日，雅化收盘 17.75 元。我们给出 12 个月目标价 17.0 元，对应约 -4.2\% 的取整后回报；未取整的多锚目标为 17.07 元，对应 -3.8\%。H1 预告确认了利润拐点，却没有确认现金转换、锂盐单位成本或 H2 持续性。现价已经接近我们的基准价值，且高于概率加权内在价值 16.95 元；在正式 H1 验证前，风险回报不支持新增仓位。
\end{houseviewbox}

\section{投资结论：利润弹性不等于可持续盈利}

雅化的关键矛盾有三层。第一，H1 预告中的利润改善是事实，Q2 归母为 7.61--9.61 亿元，说明锂盐价格、销量、销售均价和经营效率同步改善。第二，2025 年及 2026Q1 的经营现金流均为负，市场有理由要求更强的现金和营运资本验证。第三，资源自供、客户份额、实际 ASP、吨成本和民爆项目回款尚未披露，不能被提前计入基准 EPS。

\begin{exhibitbox}[表 1-1：IC 决策卡]
\begin{tabularx}{\exhibitboxwidth}{L{3.3cm}R{3.0cm}X}
\toprule
项目 & 结论 & 投资含义 \\
\midrule
当前价 / 市值 & 17.75 元 / 204.58 亿元 & 2026-07-23 收盘；总股本 11.526 亿股 \\
12 个月目标 / 行动 & 17.0 元 / 减持、避免新增 & 取整回报 -4.2\%；等待 H1 验证而非追高 \\
概率加权内在价值 & 16.95 元 & 30\% 熊、50\% 基准、20\% 牛的加权结果 \\
2026E 归母 / EPS（基准） & 20.1 亿元 / 1.74 元 & 高于 2025A，但低于三家公开卖方预测区间 \\
核心上行触发 & H1 现金、收入、毛利与营运资本共同支持 H2 & 证实单位经济后重估目标和评级 \\
核心下行触发 & 现金继续弱、应收/存货占用扩大或 H2 量价成本不支持 & 转为 watchlist only \\
\bottomrule
\end{tabularx}
\end{exhibitbox}
\sourcenote{公司《2026 年半年度业绩预告》、2025 年年报、2026Q1 报告；东财行情；AStock 情景估值模型。}

\section{盈利桥：基准 H2 不以 Q2 年化}

模型的纪律是先把 H1 预告区间纳入情景，再给 H2 设置更低、可被证伪的利润要求。基准情景以 H1 归母 12.0 亿元和 H2 归母 8.1 亿元形成 2026E 归母 20.1 亿元；这意味着即使 H1 预告中位数兑现，H2 仍需维持显著盈利，但不需要重复 Q2 中位数。熊情景对应 H2 盈利显著回落，牛情景才要求单位经济和现金同步超预期。

\begin{exhibitbox}[表 1-2：从 H1 预告到全年盈利的情景桥]
\begin{tabularx}{\exhibitboxwidth}{L{3.4cm}R{2.2cm}R{2.2cm}R{2.2cm}X}
\toprule
人民币亿元（除 EPS） & 熊 & 基准 & 牛 & 关键条件 \\
\midrule
2026H1 归母锚 & 11.0 & 12.0 & 13.0 & 均在公司预告区间内 \\
2026H2 归母假设 & 4.0 & 8.1 & 11.7 & 不把 Q2 单季简单年化 \\
2026E 归母 & 15.0 & 20.1 & 24.7 & 锂盐量价成本与民爆底座共同决定 \\
2026E EPS（元） & 1.30 & 1.74 & 2.14 & 按 11.526 亿股口径 \\
情景价值（元） & 10.3 & 17.3 & 26.1 & 周期校准 PE 的基本面锚 \\
概率 / 回报 & 30\% / -42.2\% & 50\% / -2.4\% & 20\% / +46.9\% & 相对 17.75 元现价 \\
\bottomrule
\end{tabularx}
\end{exhibitbox}

\section{行动纪律：不因估值表面偏低而抢跑}

基准情景并非断言锂价下跌，而是拒绝在单位经济与现金尚未披露时把牛情景前置。组合层面，已有仓位应把正式 H1 作为复核点：若现金、应收、存货和分部毛利支持表 1-2 的基准假设，可从“减持/避免新增”重新评估为中性；若利润与现金继续背离，低 PE 应被理解为风险折价而不是加仓信号。
